% !TeX program = pdflatex
\documentclass[12pt,a4paper]{article}

\usepackage[T2A]{fontenc}
\usepackage[utf8]{inputenc}
\usepackage[english,russian]{babel}
\usepackage{geometry}
\geometry{left=22mm,right=18mm,top=18mm,bottom=19mm}
\usepackage{setspace}
\setstretch{1.12}
\usepackage{amsmath,amssymb}
\usepackage{booktabs}
\usepackage{array}
\usepackage{longtable}
\usepackage{tabularx}
\usepackage{ragged2e}
\usepackage{hyperref}
\hypersetup{colorlinks=true,linkcolor=black,urlcolor=blue}
\usepackage{enumitem}
\setlist{itemsep=0.1em,topsep=0.2em}
\usepackage{caption}
\captionsetup{font=small,labelfont=bf}
\usepackage{xcolor}
\usepackage{url}
\urlstyle{same}
\emergencystretch=3em
\sloppy

\newcommand{\Fq}{\mathbb{F}_q}
\newcommand{\Fp}{\mathbb{F}_p}
\newcommand{\Fr}{\mathbb{F}_r}
\newcommand{\Gas}{\mathrm{gas}}
\newcommand{\code}[1]{\texttt{\detokenize{#1}}}
\newcolumntype{Y}{>{\RaggedRight\arraybackslash}X}

\begin{document}

\begin{center}
    {\Large\bfseries Итоговый отчет о выполненной работе}\\[0.4em]
    {\large Исследование стоимости сопряжений на циклических кривых\\
    в EVM и перспектив дешевого folding}\\[0.4em]
    {\normalsize MNT4-753, MNT6-753 и исследовательская конструкция lollipop-305}
\end{center}

\tableofcontents
\vspace{0.4em}

\section{Постановка задачи и практическая мотивация}

Работа посвящена проверяемым вычислениям на эллиптических кривых в EVM.
Практическая мотивация связана с рекурсивными доказательствами. Если требуется
построить цепочку доказательств, недостаточно один раз проверить отдельное
утверждение. Нужно уметь эффективно доказывать корректность предыдущего шага
внутри следующего шага. На этой идее строятся folding-схемы и рекурсивные
SNARK-системы.

В существующих реализациях, включая Sonobe/CycleFold, заметная часть стоимости
может возникать не из-за пользовательского утверждения, а из-за переноса
арифметики одной кривой в поле другой кривой. Такая арифметика называется
\emph{non-native}. Элемент большого поля раскладывается на машинные слова
(\emph{limb}-ы), а каждое умножение требует ограничений на частичные
произведения, переносы и редукцию. В качестве практического ориентира в работе
используется порядок около $9$ млн constraints для Sonobe/CycleFold-like
decider. Этот показатель не является прямым apples-to-apples сравнением с
каждым отдельным фрагментом настоящей работы, но хорошо показывает масштаб
проблемы.

Для Ethereum задача дополнительно осложняется архитектурой EVM:
\begin{itemize}
    \item EVM работает с 256-битными словами;
    \item для BN254 существуют предкомпилированные операции;
    \item для MNT4-753 и MNT6-753 аналогичных предкомпилированных операций нет;
    \item элементы MNT-полей имеют размер около 753 бит и требуют трех слов EVM.
\end{itemize}

Поэтому были исследованы два фундаментальных пути:
\begin{enumerate}
    \item выполнять арифметику сопряжения непосредственно on-chain и платить
    за нее gas;
    \item переносить вычисление off-chain и платить за доказательство,
    constraints, calldata и время работы доказчика.
\end{enumerate}

Цель работы состоит не в том, чтобы заранее выбрать один путь, а в том, чтобы
реализовать и измерить основные варианты, определить границу применимости и
проверить, дает ли цикл MNT4/MNT6 или конструкция lollipop более дешевую основу
для будущего folding.

\subsection{Статус результата}

Практическая исследовательская часть завершена как всестороннее исследование
границы применимости. Работа не ограничилась написанием одного контракта:
последовательно исследованы математические преобразования, алгоритмы
арифметики, низкоуровневая реализация в EVM, прозрачные и KZG-подходы к
проверке трассы, цикл MNT4/MNT6 и альтернативное семейство кривых
lollipop-305.

Получены следующие теоретические результаты:
\begin{itemize}
    \item формализована корректность удаления вертикальных знаменателей
    в функции Миллера для выбранной MNT4-постановки;
    \item отдельно формализована допустимость использования ненормализованных
    projective/prepared коэффициентов линий;
    \item выведено отношение со свидетельством $c$, которое заменяет длинную
    финальную экспоненту при проверке уравнения сопряжений;
    \item для MNT6 выведен отдельный знак степени:
    $r_{\mathrm{MNT6}}=q_{\mathrm{MNT6}}-N$, поэтому проверяется
    $F\cdot c^{N-q}=1$;
    \item построена формальная модель проверки цикла Миллера через
    полиномиальные отношения, KZG-opening и Merkle/ordinary-FRI;
    \item доказана корректность отношений полей цикла MNT4/MNT6;
    \item для lollipop-305 исправлена башня расширения, проверены twist- и
    Frobenius-отношения и показано, почему для третьей кривой требуется
    отдельная работа со знаменателями.
\end{itemize}

Получены следующие инженерные результаты:
\begin{itemize}
    \item реализована оптимизированная 3-limb арифметика MNT4-753 и MNT6-753
    в Solidity/Yul, включая поля расширений;
    \item реализованы и измерены Montgomery/CIOS, Barrett, FIOS, Comba/SOS,
    lazy-reduction и branchless-варианты;
    \item построены наивная Tate-модель, полное on-chain MNT4-сопряжение,
    prepared fixed-$Q$ путь и Article640 residue verifier;
    \item реализован полный MNT6 Article640 verifier с общим multi-Miller
    аккумулятором;
    \item реализованы KZG-opening over BN254, Merkle-opening слой,
    исполняемый отрицательный DEEP-FRI эксперимент и ordinary-FRI cost model;
    \item реализована 2-limb арифметика и три pairing-компонента
    исследовательской конструкции lollipop-305;
    \item добавлены cross-check-и против arkworks/Rust reference и
    воспроизводимые gas-бенчмарки.
\end{itemize}

Главный результат имеет граничный характер. Низкоуровневые оптимизации
сокращают стоимость на порядок, но без новых системных возможностей нельзя
одновременно получить малую стоимость on-chain и малую стоимость off-chain
proof-слоя простым переносом арифметики между слоями. При этом
MNT4/MNT6 cycle-native accounting показывает перспективное продолжение:
арифметику можно формулировать в родных полях цикла и тем самым избежать
наиболее дорогой BN254 non-native эмуляции.

\section{Теоретическая база вычисления сопряжения}

\subsection{Цикл Миллера и финальная экспонента}

Для reduced Tate pairing используется вычисление функции Миллера и финальная
экспонента:
\[
    e(P,Q)=f_{r,P}(Q)^{(q^k-1)/r}.
\]
Здесь $r$ --- порядок подгруппы, $q$ --- характеристика базового поля, а $k$ ---
степень вложения. Алгоритм Миллера строит $f_{r,P}$ через последовательность
удвоений и сложений. Для каждого шага обновляется аккумулятор:
\[
    f \leftarrow f^2\cdot \ell_{\mathrm{dbl}}(Q),
\]
а для ненулевой цифры signed-представления дополнительно:
\[
    f \leftarrow f\cdot \ell_{\mathrm{add}}(Q).
\]

Для уравнения сопряжений
\[
    e(P,Q)=e(R,S)
\]
удобно проверять эквивалентную форму
\[
    e(P,Q)\cdot e(-R,S)=1.
\]
Обе пары можно обрабатывать общим аккумулятором:
\[
    f \leftarrow
    f^2\cdot \ell_{Q,i}(P)\cdot \ell_{S,i}(-R).
\]
Главное преимущество общего аккумулятора: возведение в квадрат выполняется
один раз на раунд, а не отдельно для каждой пары.

\subsection{Удаление знаменателя}

В общей формуле функции Миллера возникает отношение линии к вертикальной
линии:
\[
    g_{A,B}(X)=\frac{\ell_{A,B}(X)}{v_{A+B}(X)}.
\]
Для MNT4-753 знаменатель после выбора групп лежит в собственном подполе.
Финальная экспонента уничтожает этот множитель. Поэтому в вычислительном пути
можно использовать только числитель линии. В полном тексте диплома приведено
формальное доказательство корректности такого удаления.

Практический смысл оптимизации существенен: на каждом шаге не требуется
вычислять дополнительное обращение или вести второй аккумулятор знаменателя.
Это свойство не является автоматически верным для любой кривой. Ниже будет
показано, что для третьей кривой lollipop-305 оно не выполняется в нужной
форме, и это напрямую увеличивает стоимость.

\paragraph{Два разных уровня удаления знаменателей.}
В работе используются две связанные, но не тождественные оптимизации.

Первая оптимизация относится к самой функции Миллера. В точной формуле нужно
умножать аккумулятор на
\[
    \frac{\ell_{A,B}(X)}{v_{A+B}(X)}.
\]
Если вертикальная линия $v_{A+B}(X)$ после подстановки точки лежит в
собственном подполе, ее вклад уничтожается финальной экспонентой. Поэтому
динамический цикл может не накапливать знаменатель.

Вторая оптимизация относится к построению подготовленного кэша линий.
При вычислении линии в проективных координатах ее коэффициенты определены с
точностью до ненулевого масштабирующего множителя:
\[
    (\lambda_0,\lambda_1,\lambda_2)
    \sim
    (s\lambda_0,s\lambda_1,s\lambda_2).
\]
Приведение коэффициентов к аффинной нормальной форме потребовало бы обращений
в поле. В fixed-cache пути хранятся projective/prepared коэффициенты без этой
нормализации. Масштабирующие множители допустимы только при выполнении
предпосылок конкретного pairing-equation пути: они либо уничтожаются
финальной экспонентой, либо учитываются при построении зафиксированного
кэша. Таким образом, первая оптимизация удаляет вертикальные знаменатели
Miller-функции, а вторая устраняет дорогое нормализующее деление при
подготовке line-cache.

\subsection{Проверка финальной экспоненты через свидетельство}

Статья \emph{On Proving Pairings}~\cite{eprint640} позволяет заменить полную
финальную экспоненту проверкой отношения. Для уравнения сопряжений пусть $F$ ---
результат объединенного цикла Миллера. Требуется проверить:
\[
    F^{(q^k-1)/r}=1.
\]
При выполнении условий теоремы статьи это равносильно существованию элемента
$c$, для которого:
\[
    F=c^r.
\]
Вместо длинной цепочки возведений в степень контракт проверяет эквивалентное
отношение:
\[
    F\cdot c^{-r}=1.
\]
Часть степени $c^{-r}$ встраивается в тот же signed double-and-add цикл, а
возведение в степень $q$ выполняется отображением Фробениуса.

Для MNT4 и MNT6 знаки различаются:
\[
    r_{\mathrm{MNT4}}=q_{\mathrm{MNT4}}+N,
    \qquad
    r_{\mathrm{MNT6}}=q_{\mathrm{MNT6}}-N.
\]
Следовательно:
\[
    c^{-r_{\mathrm{MNT4}}}=c^{-N-q},
    \qquad
    c^{-r_{\mathrm{MNT6}}}=c^{N-q}.
\]
Это различие было учтено в финальной реализации MNT6.

\section{Оптимизированная 3-limb арифметика MNT4-753}

\subsection{Представление данных}

Элемент базового поля MNT4-753 не помещается в одно слово EVM и хранится как
три 256-битных слова:
\[
    x=x_0+2^{256}x_1+2^{512}x_2.
\]
Для длинных цепочек умножений используется Montgomery-представление. Пусть
\[
    B=2^{256},\qquad R=B^3=2^{768}.
\]
Тогда:
\[
    \widetilde{x}=xR\bmod p,
\]
\[
    \operatorname{MontMul}(\widetilde{x},\widetilde{y})
    =\widetilde{x}\widetilde{y}R^{-1}\bmod p
    =xyR\bmod p.
\]
Результат остается в Montgomery-домене. Это важно для цикла Миллера: вход в
домен выполняется один раз, после чего тысячи операций не требуют повторного
преобразования.

\subsection{Почему выбран Montgomery/CIOS}

Были реализованы и измерены несколько вариантов:
\begin{itemize}
    \item Montgomery/CIOS;
    \item Barrett-редукция;
    \item FIOS;
    \item Comba/SOS;
    \item square-specialized Comba;
    \item branchless reduction;
    \item lazy reduction variants.
\end{itemize}

CIOS (\emph{Coarsely Integrated Operand Scanning}) объединяет накопление
частичных произведений и Montgomery-редукцию. Для EVM это выгоднее, чем
материализация полного 6-словного произведения в памяти.

\begin{table}[htbp]
\centering
\caption{Сравнение базовых 3-limb операций}
\small
\begin{tabularx}{\linewidth}{Y l r Y}
\toprule
Вариант & Операция & Gas/op & Вывод \\
\midrule
Montgomery/CIOS & $F_q$ mul & 2,959 & основной путь \\
Montgomery/CIOS & $F_q$ square & 2,947 & основной путь \\
Barrett & $F_q$ mul & 46,998 & существенно дороже \\
Barrett & $F_q$ square & 47,172 & существенно дороже \\
FIOS & $F_q$ mul & 18,509 & дороже CIOS \\
FIOS & $F_q$ square & 18,543 & дороже CIOS \\
Comba/SOS & $F_q$ mul & 18,301 & проигрывает из-за memory traffic \\
Square-Comba & $F_q$ square & 18,420 & лучше общего Comba, но хуже CIOS \\
\bottomrule
\end{tabularx}
\end{table}

\subsection{Нижняя оценка для 3-limb умножения}

Для произведения двух трехсловных чисел требуется не менее девяти попарных
произведений:
\[
    ab=\sum_{i=0}^{2}\sum_{j=0}^{2}a_i b_j 2^{256(i+j)}.
\]
Montgomery-редукция выполняет три шага. На каждом шаге добавляется произведение
корректирующего коэффициента на трехсловный модуль, то есть еще три произведения
на шаг. Получаем:
\[
    9\text{ произведений для }ab+
    9\text{ произведений для REDC}=18
\]
полных 512-битных limb-произведений.

В EVM нет инструкции \code{MULHI}. Для восстановления старшей половины
512-битного произведения используется связка:
\begin{verbatim}
lo := mul(u, v)
mm := mulmod(u, v, not(0))
hi := sub(sub(mm, lo), lt(mm, lo))
\end{verbatim}

При действующем расписании EVM опкоды имеют стоимость:
\[
    \operatorname{MUL}=5,\quad
    \operatorname{MULMOD}=8,\quad
    \operatorname{ADD}=\operatorname{SUB}=\operatorname{LT}=3.
\]
Следовательно, одно восстановление полного 512-битного произведения требует
как минимум:
\[
    5+8+3+3+3=22\ \Gas.
\]
Только неизбежные $18$ произведений дают:
\[
    18\cdot22=396\ \Gas.
\]
Это строгая нижняя граница только для восстановления произведений. Ее нельзя
сравнивать с измеренным значением как с оценкой полной функции: в ней еще нет
накопления частичных сумм, вычисления коэффициентов CIOS, нормализации и
служебных инструкций стековой машины.

Получим более содержательную нижнюю границу для арифметического каркаса
функции. Каждый 512-битный результат имеет младшую и старшую половины. Даже
если предположить идеальное расписание, обе половины необходимо добавить в
многословный аккумулятор и проверить перенос:
\[
    18\cdot 2\cdot
    (\operatorname{ADD}+\operatorname{LT})
    =
    18\cdot2\cdot(3+3)
    =
    216\ \Gas.
\]
Три CIOS-шага дополнительно вычисляют младший корректирующий коэффициент:
\[
    m_i=t_0\cdot(-q^{-1})\bmod 2^{256},
\]
то есть требуют еще трех инструкций \code{MUL}. Уже до условного финального
вычитания модуля получаем:
\[
    396+216+3\cdot5=627\ \Gas.
\]

Граница $627$ gas остается намеренно заниженной. В реальном accumulator-е
нельзя ограничиться одной проверкой переноса на каждую половину: перенос
может последовательно затрагивать следующие слова. Кроме того, необходимо
сдвигать окно CIOS после каждой итерации и в конце условно вычесть модуль.

Для проверки масштаба накладных расходов отдельно разобран фактически
исполняемый all-stack Yul-код. В нем полностью развернуты три CIOS-итерации:
циклов по limb-ам и промежуточных массивов в памяти нет. После раскрытия
вспомогательных функций структура кода содержит операции из
табл.~\ref{tab:montmul3-opcode-structure}.

\begin{table}[htbp]
\centering
\caption{Структурный opcode-состав all-stack \code{montMul3}}
\label{tab:montmul3-opcode-structure}
\small
\begin{tabular}{l r r r}
\toprule
Группа & Число & Gas/opcode & Сумма, gas \\
\midrule
\code{MUL} для limb-произведений & 18 & 5 & 90 \\
\code{MULMOD} для старших половин & 18 & 8 & 144 \\
\code{SUB} и \code{LT} внутри \code{mul512} & $36+18$ & 3 & 162 \\
\code{MUL} для коэффициентов CIOS & 3 & 5 & 15 \\
\code{ADD} накопления и переносов & до 83 & 3 & до 249 \\
\code{LT} сети переносов и финального сравнения & до 48 & 3 & до 144 \\
\code{JUMPI}/\code{JUMPDEST} условных ветвей & до 23 & $10+1$ & до 253 \\
\bottomrule
\end{tabular}
\end{table}

Таблица не является точным трассировочным отчетом для одного набора входов:
часть ветвей исполняется только при возникновении переноса, а финальное
вычитание зависит от результата. Но она показывает, почему сравнение
$2{,}959$ с числом $627$ вводило бы в заблуждение. После неизбежных
произведений значимую стоимость имеют сеть переносов и управление ветвями.
Дополнительно EVM исполняет загрузку констант, \code{PUSH}, \code{DUP},
\code{SWAP}, финальную нормализацию и служебный код, который не отражен в
арифметическом каркасе.

Следовательно, корректно фиксировать два разных результата:
\begin{itemize}
    \item $627$ gas --- строгая, но заведомо неполная нижняя граница
    арифметического каркаса;
    \item $2{,}959$ gas/op --- измеренная стоимость лучшего реализованного
    внутреннего all-stack Montgomery/CIOS пути.
\end{itemize}
Глобальная оптимальность среди всех возможных EVM-bytecode программ не
утверждается. Практическая обоснованность выбора подтверждается сравнением с
реализованными альтернативами: Barrett, FIOS, Comba/SOS, branchless reduction
и вариантами финального вычитания. Ни одна из них не снизила стоимость
минимум на 10\%.

\subsection{Башня расширений и Karatsuba}

Для MNT4 используется башня:
\[
    \mathbb{F}_{q^2}=\Fq[u]/(u^2-13),
    \qquad
    \mathbb{F}_{q^4}=\mathbb{F}_{q^2}[v]/(v^2-u).
\]
Для квадратичного расширения Karatsuba заменяет четыре умножения тремя:
\[
\begin{aligned}
    t_0&=a_0b_0,\\
    t_1&=a_1b_1,\\
    t_2&=(a_0+a_1)(b_0+b_1),\\
    c_0&=t_0+13t_1,\\
    c_1&=t_2-t_0-t_1.
\end{aligned}
\]
Та же идея применяется поверх $\mathbb{F}_{q^2}$ для $\mathbb{F}_{q^4}$.

Умножения на нерезиденты реализованы отдельными формулами. Например:
\[
    (x_0+x_1u)u=13x_1+x_0u.
\]
Это значительно дешевле общего умножения.

\begin{table}[htbp]
\centering
\caption{Стоимость операций расширений MNT4-753}
\small
\begin{tabular}{l r}
\toprule
Операция & Gas/op \\
\midrule
$\Fq$ mul & 2,959 \\
$\Fq$ square & 2,947 \\
$\mathbb{F}_{q^2}$ mul & 11,877 \\
$\mathbb{F}_{q^2}$ square & 10,592 \\
$\mathbb{F}_{q^4}$ mul & 42,764 \\
$\mathbb{F}_{q^4}$ square & 36,606 \\
\bottomrule
\end{tabular}
\end{table}

\begin{table}[htbp]
\centering
\caption{Эффект специализированных умножений на нерезиденты}
\small
\begin{tabular}{l r r r}
\toprule
Операция & Общий путь & Специализированный путь & Выигрыш \\
\midrule
умножение на $13$ в $\Fq$ & 2,961 & 1,481 & $2.0\times$ \\
умножение на $u$ в $\mathbb{F}_{q^2}$ & 12,587 & 1,731 & $7.3\times$ \\
умножение на $v$ в $\mathbb{F}_{q^4}$ & 42,253 & 1,859 & $22.7\times$ \\
\bottomrule
\end{tabular}
\end{table}

\subsection{Подготовленные разреженные линии}

После подстановки точки линия Миллера имеет специальную структуру. Ее не нужно
представлять как произвольный плотный элемент $\mathbb{F}_{q^4}$. В проекте
реализованы:
\begin{itemize}
    \item prepared fixed-$Q$ cache;
    \item sparse-формат коэффициентов;
    \item специализированное умножение аккумулятора на линию;
    \item pointer/packed memory API;
    \item потоковое чтение code-shards через \code{EXTCODECOPY};
    \item экспериментальное слияние $f\leftarrow f^2\ell(P)$ без лишнего
    промежуточного объекта.
\end{itemize}

Последний агрессивный fused-вариант корректен, но дает лишь около $20$ тыс.
gas экономии на полном single-pairing пути. Это показывает, что после
оптимизации памяти доминирует уже не копирование данных, а собственно
арифметика расширений.

\paragraph{Формальная идея разреженного умножения.}
Пусть аккумулятор цикла Миллера имеет общий вид:
\[
    f=f_0+f_1v,\qquad f_0,f_1\in\mathbb{F}_{q^2}.
\]
После подстановки $G_1$-точки подготовленная линия не является произвольным
элементом $\mathbb{F}_{q^4}$: часть ее коэффициентов равна нулю или выражается
через заранее известные координаты. Условно ее можно записать как:
\[
    \ell(P)=\ell_0(P)+\ell_1(P)v,
\]
где $\ell_0,\ell_1$ вычисляются из компактного набора prepared-коэффициентов.
Вместо построения плотного объекта и вызова общего умножения
\[
    f\leftarrow\operatorname{Mul}_{\mathbb{F}_{q^4}}(f,\ell(P))
\]
используется специализированная формула:
\[
\begin{aligned}
    t_0&=f_0\ell_0,\\
    t_1&=f_1\ell_1,\\
    t_2&=(f_0+f_1)(\ell_0+\ell_1),\\
    f'_0&=t_0+u\,t_1,\\
    f'_1&=t_2-t_0-t_1.
\end{aligned}
\]
Умножение на $u$ выполняется дешевым переставлением коэффициентов. В fused
hot path вычисление
\[
    f\leftarrow f^2\ell(P)
\]
дополнительно сливается с загрузкой коэффициентов: промежуточный плотный
$\mathbb{F}_{q^4}$-объект не материализуется.

\subsection{Низкоуровневые оптимизации EVM}

Алгоритмические формулы были дополнены низкоуровневой реализацией:
\begin{itemize}
    \item фиксированный размер $3\times3$ limb развернут вручную, без
    динамических циклов и динамических массивов;
    \item старшая половина произведения восстанавливается через
    \code{MUL}/\code{MULMOD}, так как EVM не имеет \code{MULHI};
    \item CIOS объединяет умножение и REDC, чтобы не сохранять полное
    шестисловное произведение;
    \item промежуточные значения по возможности удерживаются на стеке;
    \item pointer/packed API и scratch arena уменьшают копирование вложенных
    структур расширений;
    \item lazy reduction применяется только там, где диапазон промежуточного
    значения формально ограничен и отложенная редукция действительно дешевле;
    \item специальные операции \code{mulBy13}, \code{mulByU},
    \code{mulByV} заменяют общие умножения;
    \item prepared-кэш может читаться потоково из data-контрактов через
    \code{EXTCODECOPY}, без повторной передачи большого blob;
    \item ветвистая и branchless-редукции реализованы рядом и измерены;
    branchless-вариант оставлен исследовательским, поскольку в EVM он дороже.
\end{itemize}

Эта часть закрывает вопрос не только о выборе математического алгоритма, но и
о том, какие расходы остаются после ручной адаптации к стековой архитектуре
EVM.

\section{MNT4: от наивного Tate pairing к оптимизированному пути}

\subsection{Наивная точка отсчета}

Для честного сравнения построена наивная модель Tate pairing. Она намеренно не
использует:
\begin{itemize}
    \item короткий Ate loop;
    \item prepared lines;
    \item sparse multiplication;
    \item Yul hot path;
    \item Frobenius-декомпозицию финальной экспоненты.
\end{itemize}

Полный неоптимизированный вызов превышает практические лимиты EVM, поэтому
измерены корректные микроблоки и построена строгая нижняя экстраполяция:
\[
\begin{aligned}
    \text{Miller} &=
    752\cdot 973{,}261+372\cdot 486{,}819\\
    &=912{,}988{,}940\ \Gas,\\
    \text{FE} &=
    2259\cdot 486{,}899+1100\cdot 486{,}819\\
    &=1{,}635{,}405{,}741\ \Gas.
\end{aligned}
\]
Итого:
\[
    2{,}548{,}394{,}681\ \Gas.
\]
Это нижняя оценка: она не включает построение линий, проверки входов и часть
накладных расходов.

\subsection{Оптимизационная лестница}

Не все оптимизации можно честно представить как последовательность полных
контрактов: некоторые меняют стоимость одной операции поля, другие ---
длину цикла Миллера, третьи --- способ доставки кэша. Поэтому ниже приведены
три взаимосвязанные таблицы.

\begin{table}[htbp]
\centering
\caption{Интегральная лестница полного MNT4-пути}
\small
\begin{tabularx}{\linewidth}{r Y r r}
\toprule
Шаг & Реализация & Gas & Снижение к предыдущему шагу \\
\midrule
0 & Наивная Tate cost model & 2,548,394,681 & --- \\
1 & Полный optimized fixed-$Q$ on-chain путь & 258,753,182 & 89.85\% \\
2 & Fixed-$Q$ prepared sparse blob & 79,588,799 & 69.24\% \\
3 & Fixed-$Q$ prepared sparse code-shards & 79,586,596 & 0.003\% \\
\bottomrule
\end{tabularx}
\end{table}

\begin{table}[htbp]
\centering
\caption{Локальные оптимизации арифметики}
\small
\begin{tabularx}{\linewidth}{Y Y r r}
\toprule
Оптимизация & Что меняется & Было & Стало \\
\midrule
Montgomery/CIOS + Yul & $F_q$ mul вместо generic Solidity multi-limb &
18,770 & 2,959 \\
Karatsuba для башни & $F_{q^4}$ mul вместо generic расширения &
486,819 & 42,764 \\
Умножение на $13$ & общая операция заменяется сложениями &
2,961 & 1,481 \\
Умножение на $u$ & $(x_0+x_1u)u=13x_1+x_0u$ &
12,587 & 1,731 \\
Умножение на $v$ & перестановка $F_{q^2}$-коэффициентов &
42,253 & 1,859 \\
\bottomrule
\end{tabularx}
\end{table}

\begin{table}[htbp]
\centering
\caption{Лестница цикла Миллера, финальной экспоненты и кэша}
\small
\begin{tabularx}{\linewidth}{Y Y r r}
\toprule
Оптимизация & Смысл & Было & Стало \\
\midrule
Ate вместо полного Tate loop & короткая signed-запись параметра &
912,988,940 & 427,284,953 \\
Sparse line multiplication & линия не материализуется как общий $F_{q^4}$ &
973,261/шаг & $\approx137,874$/digit \\
Prepared fixed-$Q$ cache & line generation переносится из on-chain &
258,753,182 & 79,588,799 \\
Frobenius/$w_0$ FE & декомпозиция вместо binary exponentiation &
1,635,405,741 & $\approx28,134,825$ \\
Aggressive fused hot path & слияние $f^2$ и sparse line multiplication &
51,978,344 & 51,949,399 \\
Code-shards & повторный blob заменяется \code{EXTCODECOPY} &
79,588,799 & 79,586,596 \\
\bottomrule
\end{tabularx}
\end{table}

Для строки Ate loop прочие условия специально оставлены неизменными:
\[
    752\cdot973{,}261+372\cdot486{,}819
    =912{,}988{,}940\ \Gas,
\]
\[
    377\cdot973{,}261+124\cdot486{,}819
    =427{,}284{,}953\ \Gas.
\]
Так отдельно виден эффект сокращения параметра цикла. Для финальной
экспоненты измеренный overhead optimized пути вычисляется как:
\[
    80{,}113{,}169-51{,}978{,}344
    =28{,}134{,}825\ \Gas.
\]

Оптимизации дают более чем десятикратное снижение относительно наивного пути.
Однако даже prepared-вариант остается дороже обычного лимита блока Ethereum.
После применения fused hot path дальнейшие локальные улучшения дают уже доли
процента: доминирует число неизбежных операций в расширенном поле.

\subsection{Формальная нижняя оценка стоимости Miller core}

Для двухпарного уравнения:
\[
    e(P,Q)e(-R,S)=1
\]
используем:
\[
    L=376,\qquad H=123,
\]
где $L$ --- число удвоений, $H$ --- число ненулевых addition-шагов.
Минимально нужны:
\begin{itemize}
    \item $L$ возведений аккумулятора в квадрат;
    \item $2L$ применений doubling-линий;
    \item $2H$ применений addition-линий.
\end{itemize}
Грубая нижняя оценка через измеренные $\mathbb{F}_{q^4}$-операции:
\[
\begin{aligned}
    376\cdot 36{,}606 &= 13{,}763{,}856,\\
    2\cdot376\cdot42{,}764 &= 32{,}157{,}728,\\
    2\cdot123\cdot42{,}764 &= 10{,}499{,}864.
\end{aligned}
\]
Суммарно:
\[
    56{,}421{,}448\ \Gas.
\]
Реальный код дополнительно загружает sparse-коэффициенты, вычисляет значения
линий, работает с памятью и обрабатывает signed loop. Поэтому measured
Miller core порядка $87.7$ млн gas согласуется с формальной моделью.

\section{MNT4 direct residue verifier по ePrint 2024/640}

\subsection{API и граница доверия}

В основном fixed-shards режиме контракт проверяет:
\[
    e(P,Q)e(-R,S)=1,
\]
где $Q,S\in G_2$ фиксированы при развертывании, а $P,R\in G_1$ передаются
пользователем. Prepared-кэш линий хранится в runtime-коде data-контрактов.
Адреса shards фиксируются в конструкторе. Пользователь не может заменить их
при вызове.

Это доверенная предрегистрация fixed-cache: контракт защищен от подмены кэша
во время вызова, но не доказывает заново, что развертыватель построил линии
правильно. Для произвольной меняющейся точки $Q$ нужен отдельный proof корректности
кэша или прямой пересчет.

\subsection{Результаты MNT4}

\begin{table}[htbp]
\centering
\caption{Direct residue verifier MNT4}
\small
\begin{tabularx}{\linewidth}{Y Y r}
\toprule
Сценарий & Что выполняется & Gas \\
\midrule
Miller core & цикл Миллера без финальной экспоненты & 87,747,219 \\
Full equation & Miller core + обычная FE & 115,997,313 \\
Residue equation & Miller core + $c$-проверка вместо полной FE & 93,879,746 \\
Fixed-shards residue equation с проверкой $G_1$ & основной fixed-cache API & 93,734,789 \\
\bottomrule
\end{tabularx}
\end{table}

Вклад обычной FE:
\[
    115{,}997{,}313-87{,}747{,}219
    =28{,}250{,}094\ \Gas.
\]
Вклад residue-проверки:
\[
    93{,}879{,}746-87{,}747{,}219
    =6{,}132{,}527\ \Gas.
\]
Экономия:
\[
    22{,}117{,}567\ \Gas.
\]
Оптимизация финальной экспоненты работает, но не устраняет главный расход:
цикл Миллера.

\section{Проверка цикла Миллера через полиномиальные отношения}

\subsection{Математическая идея}

Следующий шаг статьи~\cite{eprint640} --- не исполнять все умножения
расширенного поля on-chain, а проверить их как полиномиальные отношения.
В статье сформулирована общая идея. Для MNT4-753 потребовалась отдельная
адаптация: нужно связать quotient-check не с абстрактными умножениями, а с
реальными состояниями общего multi-Miller аккумулятора, fixed-$Q$/fixed-$S$
кэшем и residue-свидетельством $c$.

Пусть:
\[
    \mathbb{F}_{q^4}\simeq \Fq[X]/(p(X)).
\]
Для умножения $a\cdot b=c$ существует частный многочлен $Q(X)$:
\[
    a(X)b(X)-c(X)=Q(X)p(X).
\]
Verifier выбирает случайную точку $\alpha$ и проверяет:
\[
    a(\alpha)b(\alpha)-c(\alpha)=Q(\alpha)p(\alpha).
\]
Для набора отношений используется challenge $\beta$:
\[
    \sum_i\beta^i R_i(X)=Q(X)Z_H(X).
\]
Для цикла Миллера $R_i$ кодируют реальные переходы:
\[
    f_{i+1}=f_i^2\ell_{Q,i}(P)\ell_{S,i}(-R).
\]

Чтобы prover не подменил значения после получения $\alpha$, нужны commitments
и доказательства открытия.

\subsection{Построенная MNT4-модель}

Для MNT4 была построена следующая модель.

\paragraph{Блочное сжатие цикла.}
Несколько раундов объединяются в один блок. Для двух шагов:
\[
    f_{i+2}
    =
    f_i^4\ell_i^2\ell_{i+1}.
\]
Для блока длины $d$:
\[
    f_{i+d}
    =
    f_i^{2^d}
    \prod_{j=0}^{d-1}\ell_{i+j}^{2^{d-1-j}}.
\]
В модели используется $d=5$. Тогда один блочный переход для уравнения
сопряжений имеет вид:
\[
    F_{b+1}
    =
    F_b^{32}
    G_b^Q(P)
    G_b^S(-R)
    c^{\kappa_b}.
\]
Здесь $G_b^Q$ и $G_b^S$ --- сжатые делители для зафиксированных точек
$Q,S$, а $\kappa_b$ --- заранее известный показатель residue-рекурсии.

\paragraph{Проверка вычисления делителей.}
Для fixed-$Q$/fixed-$S$ коэффициенты делителей строятся off-chain и
фиксируются root-ом. Значение $G_b^Q(P)$ не должно выбираться prover-ом
произвольно. Поэтому строится таблица вычисления по схеме Горнера:
\[
    h_{j+1}=h_j P_x+\gamma_j,
\]
где $\gamma_j$ --- зафиксированный коэффициент сжатого делителя. Аналогично
вычисляется $G_b^S(-R)$. Локальные отношения Горнера и блочные переходы
собираются в quotient polynomial.

\paragraph{Что доказано для модели.}
Для спецификации были формализованы:
\begin{itemize}
    \item связь пяти пошаговых переходов Миллера с одним сжатым делителем;
    \item корректность Horner-рекурсии для fixed-divisor коэффициентов;
    \item quotient identity для локальных переходов;
    \item Merkle-binding динамических таблиц;
    \item Fiat--Shamir transcript и детерминированный порядок запросов;
    \item ordinary-FRI low-degree проверка;
    \item итоговая soundness-оценка как сумма ошибок FRI,
    Schwartz--Zippel, Merkle-binding и Fiat--Shamir модели.
\end{itemize}

Таким образом, исследован не эвристический ``кэш строк'', а конкретная
неинтерактивная схема проверки реальных переходов цикла Миллера.

\subsection{KZG over BN254}

KZG дает короткое доказательство открытия и дешевую on-chain проверку благодаря
BN254 precompile. В проекте реализован реальный opening layer:
\[
    C_j=\operatorname{Commit}(w_j),\qquad
    C_Q=\operatorname{Commit}(Q).
\]
Для одного полинома проверяется стандартное отношение:
\[
    e(C-yG_1+x\pi,G_2)=e(\pi,\tau G_2).
\]
Здесь $C$ --- commitment, $x$ --- точка открытия, $y=f(x)$ ---
заявленное значение, $\pi$ --- opening proof, $\tau G_2$ --- элемент SRS.
Measured стоимость одной on-chain opening-проверки составляет около:
\[
    133{,}039\ \Gas.
\]

Проблема переносится внутрь circuit: MNT4-арифметика становится non-native
относительно BN254.

\begin{table}[htbp]
\centering
\caption{Оценка BN254 non-native constraints}
\small
\begin{tabular}{l r}
\toprule
Операция & Constraints \\
\midrule
MNT4 $\Fq$ mul как non-native & 3,692 \\
MNT4 $\mathbb{F}_{q^2}$ mul как non-native & 11,300 \\
Модель MNT4 $\mathbb{F}_{q^4}$ mul как non-native & 63,436 \\
Приближенная sparse Miller relation & 63,309,128 \\
\bottomrule
\end{tabular}
\end{table}

Число $63{,}309{,}128$ получается не из размера KZG proof. Оно относится к
relation, которую нужно доказать до открытия: каждое MNT4 $F_{q^4}$
умножение при переносе в BN254 R1CS раскладывается на non-native операции.
В проекте отдельно измерены $3{,}692$ constraints для одного MNT4
$F_q$-умножения через emulated field, $11{,}300$ для $F_{q^2}$ и
$63{,}436$ для модели $F_{q^4}$. Подстановка числа sparse Miller
переходов дает ориентир $63.3$ млн constraints. То есть KZG уменьшает
gas verifier-а, но возвращает исходную проблему роста off-chain constraints.

\subsection{Merkle/FRI}

В альтернативном направлении prover коммитится к таблицам через Merkle root.
Verifier проверяет Merkle-пути открытых строк. FRI
(\emph{Fast Reed--Solomon Interactive Oracle Proof of Proximity}) доказывает,
что открываемые таблицы согласуются с многочленами ограниченной степени.
Fiat--Shamir-преобразование делает схему неинтерактивной: challenges
детерминированно вычисляются из transcript.

В проекте реализованы:
\begin{itemize}
    \item Merkle opening layer;
    \item исполняемый архивный DEEP-FRI отрицательный эксперимент;
    \item спецификация блочно-сжатого ordinary-FRI пути;
    \item актуальная ordinary-FRI cost model для нативного MNT4-поля.
\end{itemize}
Полноценная production-замена Miller loop не заявляется: результат используется
как строгая stop/go модель стоимости.

Для MNT4 один элемент поля занимает:
\[
    3\cdot 32=96\text{ байт}.
\]
Умеренная модель открытий:
\[
    4096\cdot96+128\cdot16\cdot32+16\cdot32
    =459{,}264\text{ байт}.
\]
Актуальная ordinary-FRI модель для 128-битного профиля:

\begin{table}[htbp]
\centering
\caption{Ordinary-FRI cost model}
\small
\begin{tabular}{l r}
\toprule
Показатель & Значение \\
\midrule
Число запросов & 576 \\
Оценка soundness & 135.76 бит \\
FRI schedule & $[2,2,2,4]$ \\
Размер доказательства & 2,072,224 байт \\
Строгая нижняя оценка & 51,359,352 gas \\
Ожидаемая модельная оценка & 78,340,624 gas \\
\bottomrule
\end{tabular}
\end{table}

\paragraph{Из чего складывается нижняя оценка.}
Для каждого случайного запроса verifier должен получить раскрытые значения
столбцов, соседние значения для локального перехода, Merkle-path и элементы
слоев FRI. Поскольку один MNT4-элемент занимает $96$ байт, даже до исполнения
арифметики calldata становится существенным компонентом стоимости. Актуальная
модель суммирует:
\[
    \Gas_{\min}
    =
    \Gas_{\mathrm{calldata}}
    +
    \Gas_{\mathrm{Merkle}}
    +
    \Gas_{\mathrm{FRI\ folding}}
    +
    \Gas_{\mathrm{local\ checks}}.
\]
Для профиля с $576$ запросами строгая нижняя оценка равна
$51{,}359{,}352$ gas. Она учитывает минимально необходимые данные и операции,
но не претендует на стоимость готового Solidity-контракта. Исполняемая модель
с реалистичными накладными расходами дает:
\[
    78{,}340{,}624\ \Gas.
\]

\paragraph{Почему не выбран DEEP-FRI production-путь.}
Более консервативная DEEP-FRI микротрасса была реализована как отрицательный
эксперимент. Она проверяет более широкую временную трассу и требует большего
числа раскрытий. После профилирования этот путь перенесен в
\code{research_variants}: он криптографически содержателен, но не дает
практического gas-выигрыша.

Вывод: Merkle/ordinary-FRI может приблизиться к Article640 fixed-shards
verifier и в оптимистичной нижней модели пересечь уровень $60$ млн gas.
Однако реалистичная оценка остается около $78.3$ млн gas. Стоимость
переносится из прямых $F_{q^4}$-умножений в calldata, Merkle hashing,
раскрытия таблиц и FRI folding. Поэтому направление подробно исследовано,
но не принято как финальная production-замена цикла Миллера.

\section{MNT6-753 и цикл MNT4/MNT6}

\subsection{Почему нужна вторая кривая}

MNT4-753 и MNT6-753 образуют цикл полей:
\[
    \Fr(\mathrm{MNT4})=\Fq(\mathrm{MNT6}),
\]
\[
    \Fr(\mathrm{MNT6})=\Fq(\mathrm{MNT4}).
\]
Для будущего folding это принципиально: арифметику одной стороны можно
выражать ближе к родному полю другой стороны, избегая наиболее грубой
BN254 non-native эмуляции.

\subsection{MNT6 arithmetic и residue verifier}

Для MNT6 реализована башня:
\[
    \Fq\longrightarrow\mathbb{F}_{q^3}\longrightarrow\mathbb{F}_{q^6}.
\]
Реализованы packed/pointer API, scratch arena, prepared cache, потоковое чтение
code-shards, fused line multiplication и Frobenius/w0 FE.

\begin{table}[htbp]
\centering
\caption{Базовые операции MNT6}
\small
\begin{tabular}{l r}
\toprule
Операция & Gas/op \\
\midrule
$\Fq$ mul & 2,314 \\
$\Fq$ square & 2,301 \\
$\mathbb{F}_{q^3}$ mul & 36,472 \\
$\mathbb{F}_{q^3}$ square & 26,813 \\
$\mathbb{F}_{q^6}$ mul & 166,774 \\
$\mathbb{F}_{q^6}$ square & 113,788 \\
\bottomrule
\end{tabular}
\end{table}

После финального аудита MNT6 verifier приведен к той же архитектуре, что и
MNT4: используется общий multi-Miller аккумулятор и один квадрат на раунд.
Из $r_{\mathrm{MNT6}}=q-N$ следует:
\[
    F\cdot c^{N-q}=F\cdot c^{-r}=1.
\]

\paragraph{Лемма 1. Корректность башни расширений MNT6.}
Пусть:
\[
    \mathbb{F}_{q^3}=\Fq[v]/(v^3-11),
    \qquad
    \mathbb{F}_{q^6}=\mathbb{F}_{q^3}[w]/(w^2-v).
\]
При выборе нерезидентов из параметров MNT6-753 эти многочлены неприводимы,
поэтому конструкции задают поля степеней $3$ и $6$. Следовательно,
Karatsuba-формулы, отображения Фробениуса и обращение элементов реализуются в
том же поле, в котором лежит результат MNT6-сопряжения. Константы и
reference-значения кросс-проверены с arkworks.

\paragraph{Лемма 2. Общий аккумулятор для уравнения сопряжений.}
Пусть на раунде $i$ индивидуальные аккумуляторы обновляются как:
\[
    f'_{Q}=f_Q^2\ell_{Q,i}(P),
    \qquad
    f'_{S}=f_S^2\ell_{S,i}(-R).
\]
Для произведения $F=f_Qf_S$ имеем:
\[
    F'
    =
    f'_Qf'_S
    =
    (f_Qf_S)^2
    \ell_{Q,i}(P)\ell_{S,i}(-R).
\]
Следовательно, вместо двух возведений аккумуляторов в квадрат достаточно
одного. Именно эта формула используется в MNT6 fixed-shards verifier.

\paragraph{Теорема 1. Residue-отношение для MNT6.}
Для параметров MNT6:
\[
    r_{\mathrm{MNT6}}=q_{\mathrm{MNT6}}-N.
\]
Если результат общего цикла Миллера равен $F$ и prover предоставляет
$c$ такое, что $F=c^r$, то:
\[
    F\cdot c^{-r}
    =
    F\cdot c^{N-q}
    =
    1.
\]
Возведение в степень $q$ реализуется отображением Фробениуса. Поэтому
контракт проверяет то же математическое условие, что и полная финальная
экспонента для уравнения сопряжений, но не выполняет полный exponentiation
chain.

\begin{table}[htbp]
\centering
\caption{Результаты MNT6}
\small
\begin{tabularx}{\linewidth}{Y r}
\toprule
Режим & Gas \\
\midrule
Miller loop, packed pointer blob & 93,254,054 \\
Packed Frobenius/w0 FE & 38,428,108 \\
Полное single pairing: Miller + packed FE & 131,685,843 \\
Диагностический residue digest одного pairing & 103,277,505 \\
Два Miller loop + полная FE для equation & 226,078,963 \\
Общий multi-Miller + $c$-проверка для equation & 172,004,717 \\
\bottomrule
\end{tabularx}
\end{table}

Экономия общего residue verifier относительно контрольного equation baseline:
\[
    226{,}078{,}963-172{,}004{,}717
    =54{,}074{,}246\ \Gas,
\]
то есть $23.92\%$.

\subsection{Cycle-native accounting}

Rust-модуль проверяет параметры цикла через arkworks, строит reference
pairings и считает operation/constraints accounting для будущего native
relation layer.

\paragraph{Теорема 2. Связь полей цикла.}
Для выбранных параметров:
\[
    \Fr(\mathrm{MNT4})=\Fq(\mathrm{MNT6}),
    \qquad
    \Fr(\mathrm{MNT6})=\Fq(\mathrm{MNT4}).
\]
Равенства проверяются программно сравнением модулей arkworks. Поэтому
базовые умножения одной стороны цикла могут выражаться как нативные
multiplication constraints в скалярном поле другой стороны, а не как
753-битная арифметика, эмулируемая поверх BN254.

\begin{table}[htbp]
\centering
\caption{Предварительная модель MNT-cycle-native relation}
\small
\begin{tabular}{l l r r r}
\toprule
Сторона & Башня & Miller rounds & Additions & Prepared relation \\
\midrule
MNT4 & $\mathbb{F}_{q^2}/\mathbb{F}_{q^4}$ & 376 & 124 & 24,126 \\
MNT6 & $\mathbb{F}_{q^3}/\mathbb{F}_{q^6}$ & 376 & 123 & 48,942 \\
\midrule
\multicolumn{4}{r}{Сумма строительных блоков} & 73,068 \\
\bottomrule
\end{tabular}
\end{table}

Для сравнения:
\[
    \text{BN254 non-native sparse Miller estimate}=63{,}309{,}128,
\]
\[
    \text{Sonobe/CycleFold-like anchor}\approx9{,}000{,}000.
\]
Число $73{,}068$ нельзя называть стоимостью готового CycleFold. Это ручная
модель multiplication constraints для подготовленных relation fragments.
Она показывает порядок потенциального выигрыша нативной постановки и
обосновывает направление следующего этапа.

\paragraph{Оценка будущего cycle-native folding слоя.}
Внутреннее ядро relation уже посчитано по реальным формулам проекта:
\[
    24{,}126+48{,}942=73{,}068
\]
нативных multiplication constraints. Для MNT4 ручная модель дополнительно
сверялась с собранным R1CS-фрагментом: расхождение составляет единицы
ограничений. Полный CycleFold-circuit еще не реализован, поэтому нельзя
выдавать одно точное итоговое число. Можно честно построить чувствительность
к накладным расходам folding-обвязки:

\begin{table}[htbp]
\centering
\caption{Проекция cycle-native constraints при разных накладных расходах}
\small
\begin{tabular}{l r r}
\toprule
Сценарий & Коэффициент к relation core & Constraints \\
\midrule
Только подготовленные fragments & $1\times$ & 73,068 \\
Легкая folding-обвязка & $2\times$ & 146,136 \\
Умеренная folding-обвязка & $4\times$ & 292,272 \\
Консервативная folding-обвязка & $8\times$ & 584,544 \\
Верхняя исследовательская граница & $12\times$ & 876,816 \\
\bottomrule
\end{tabular}
\end{table}

Эта таблица не является замером готового CycleFold. Она отвечает на вопрос,
какой запас допускает измеренное ядро до уровня Sonobe/CycleFold-like
ориентира около $9$ млн constraints. Даже при двенадцатикратной обвязке
оценка остается меньше $1$ млн. Следовательно, полученные результаты
поддерживают основной тезис работы: MNT4/MNT6 цикл потенциально позволяет
кратно снизить constraints, если следующий этап реализовать нативно, а не
переносить MNT-арифметику в BN254.

Если механически сложить лучшие прямые EVM equation verifier-ы двух сторон,
получится:
\[
    93{,}734{,}789+172{,}004{,}717
    =265{,}739{,}506\ \Gas.
\]
Эта сумма иллюстрирует границу прямого on-chain исполнения. Будущий folding
не должен выполнять оба verifier-а on-chain; его смысл как раз в переносе
relation layer в нативный цикл.

\section{Исследовательское направление lollipop-305}

\subsection{Идея}

Статья \emph{Cycles of supersingular elliptic curves for pairing-based proof
systems}~\cite{eprint1627} предлагает конструкции, в которых часть
арифметики выполняется над меньшими полями. Для исследовательского примера
lollipop-305 выбран параметр:
\[
    x=8004046504391788107635887004283725454478544674,
\]
\[
    p=x^2-x+1,\qquad q=x^2+1.
\]
Оба основных поля имеют размер около 305 бит и помещаются в два слова EVM.
Цель этого направления --- проверить, является ли уменьшение числа limb-ов
достаточным условием практического выигрыша. Это не production-предложение:
параметры lollipop-305 выбраны как исследовательский прототип для измерения
архитектурного эффекта.

\subsection{Кривые конструкции}

\begin{table}[htbp]
\centering
\caption{Pairing-компоненты lollipop-305}
\small
\begin{tabularx}{\linewidth}{l Y Y}
\toprule
Часть & Кривая и поле & Целевое поле \\
\midrule
stick & pairing-friendly prototype над $\Fp$ & $\mathbb{F}_{p^4}$ \\
cycle $E$ & $E/\mathbb{F}_{p^2}: y^2=x^3+(\mu+1)x$ & $\mathbb{F}_{p^4}$ \\
cycle $\widehat{E}$ &
$\widehat{E}/\mathbb{F}_{q^2}: y^2=x^3+(\eta+2)$ &
$\mathbb{F}_{q^6}$ \\
\bottomrule
\end{tabularx}
\end{table}

Полная конструкция содержит три pairing-компонента:
\begin{enumerate}
    \item ``палочку'' --- pairing-friendly prototype над $\Fp$;
    \item первую supersingular cycle-кривую $E/\mathbb{F}_{p^2}$;
    \item вторую supersingular cycle-кривую
    $\widehat{E}/\mathbb{F}_{q^2}$.
\end{enumerate}
Дополнительно статья использует меньшую непаринговую кривую для рекурсивного
сценария, но она не является отдельным EVM pairing verifier.

Проверены отношения:
\[
    \#E(\mathbb{F}_{p^2})=p^2+1=qN_q,
\]
\[
    \#\widehat{E}(\mathbb{F}_{q^2})=q^2-q+1=pN_p.
\]

\paragraph{Лемма 3. Исправление башни для stick/$E$-части.}
Первоначальный кандидат:
\[
    \mathbb{F}_{p^4}
    =
    \mathbb{F}_{p^2}[v]/(v^2-u),
    \qquad u^2=-1,
\]
не задает нужное поле: для выбранного $p$ элемент $u$ является квадратом в
$\mathbb{F}_{p^2}$. В работе вычислительно найден и проверен неквадрат:
\[
    \xi=1+u.
\]
Поэтому используется исправленная башня:
\[
    \mathbb{F}_{p^4}
    =
    \mathbb{F}_{p^2}[v]/(v^2-\xi).
\]

\paragraph{Лемма 4. Twist/Frobenius-отношение.}
Для исправленной башни построено отображение из twist в исходную кривую:
\[
    \psi(X,Y)=(\xi X,\xi vY).
\]
Rust backend проверяет:
\[
    [r]P=\mathcal O,\quad
    [r]\psi(Q')=\mathcal O,\quad
    \pi_p(\psi(Q'))=[p]\psi(Q'),
\]
а также нетривиальность результата сопряжения и условие принадлежности
целевой подгруппе.

\paragraph{Лемма 5. Корректная форма для $\widehat E$.}
Для третьей кривой выбрана башня:
\[
    \mathbb{F}_{q^6}
    =
    \mathbb{F}_{q^2}[w]/(w^3-\rho),
\]
где:
\[
    \rho^2=\frac{2+\eta}{2-\eta}.
\]
Distortion/Frobenius map имеет вид:
\[
    \psi_{\widehat E}(x,y)=(w^2x^q,\rho y^q).
\]
Для $\widehat E$ упрощенная Tate-style формула не проходит. Корректным
reference путем является Weil-отношение:
\[
    e_p(P,Q)
    =
    (-1)^p\frac{f_{p,P}(Q)}{f_{p,Q}(P)}.
\]
Отсюда следует verifier relation:
\[
    f_{p,P}(Q)f_{p,-P}(Q)
    =
    f_{p,Q}(P)f_{p,Q}(-P).
\]

\subsection{Оптимизированная 2-limb арифметика}

Для двухсловного Montgomery-умножения generic lower bound содержит:
\[
    4\text{ произведения для }ab+
    4\text{ произведения для REDC}.
\]
Для lollipop-305 старший limb модуля меньше $2^{49}$. Поэтому:
\[
    a_1b_1<2^{98}<2^{256}.
\]
Одно полное 512-битное восстановление можно заменить обычным \code{MUL}.
Дополнительно убрана запись limb-а, который сразу удаляется CIOS-сдвигом.

Эта часть не является упрощенным демонстрационным скриптом. Для 305-битных
полей проделан тот же класс инженерной работы, что и для 753-битной
арифметики MNT4:
\begin{itemize}
    \item Montgomery/CIOS-представление;
    \item ручной Yul hot path;
    \item специальная обработка малого старшего limb-а;
    \item арифметика расширений $\mathbb{F}_{p^2}$,
    $\mathbb{F}_{p^4}$, $\mathbb{F}_{q^2}$ и $\mathbb{F}_{q^6}$;
    \item Karatsuba и специализированные умножения на нерезиденты;
    \item prepared lines, residue-пути и fixed-cache режимы;
    \item Rust cross-check и gas-бенчмарки.
\end{itemize}

\begin{table}[htbp]
\centering
\caption{Сравнение 2-limb lollipop-305 и 3-limb MNT4}
\small
\begin{tabular}{l r r r}
\toprule
Операция & lollipop-305 & MNT4-753 & Выигрыш \\
\midrule
$F_p$ mul & 1,018 & 2,959 & $2.91\times$ \\
$F_p$ square & 940 & 2,947 & $3.14\times$ \\
$F_{p^2}$ mul & 4,171 & 11,877 & $2.85\times$ \\
$F_{p^2}$ square & 3,254 & 10,592 & $3.26\times$ \\
$F_{p^4}$ mul & 14,942 & 42,764 & $2.86\times$ \\
$F_{p^4}$ square & 13,211 & 36,606 & $2.77\times$ \\
\bottomrule
\end{tabular}
\end{table}

\subsection{Почему для третьей кривой нужны два аккумулятора}

Для stick и $E$ можно использовать Article640-style residue verifier. Для
$\widehat{E}/\mathbb{F}_{q^2}$ степень вложения равна $3$. Стандартное удаление
знаменателей линий не работает в требуемой форме. Поэтому backend и Solidity
отдельно накапливают:
\[
    F_{\mathrm{num}},\qquad F_{\mathrm{den}}.
\]
На doubling-шаге:
\[
    F_{\mathrm{num}}\leftarrow
    F_{\mathrm{num}}^2\ell_{\mathrm{dbl}}(P),
\]
\[
    F_{\mathrm{den}}\leftarrow
    F_{\mathrm{den}}^2v_{\mathrm{dbl}}(P).
\]
На addition-шаге:
\[
    F_{\mathrm{num}}\leftarrow
    F_{\mathrm{num}}\ell_{\mathrm{add}}(P),
\]
\[
    F_{\mathrm{den}}\leftarrow
    F_{\mathrm{den}}v_{\mathrm{add}}(P).
\]
В конце проверяется отношение:
\[
    c^p F_{\mathrm{den}}=F_{\mathrm{num}}.
\]
Это эквивалентно residue-проверке частного:
\[
    F=F_{\mathrm{num}}/F_{\mathrm{den}}.
\]
Ведение двух аккумуляторов и 305-битное возведение $c^p$ в
$\mathbb{F}_{q^6}$ делают третью часть заметно дороже.

\subsection{Результаты lollipop-305}

\begin{table}[htbp]
\centering
\caption{Pairing-verifier-ы lollipop-305}
\small
\begin{tabularx}{\linewidth}{Y Y r}
\toprule
Часть & Режим & Function gas \\
\midrule
stick & residue FE & 8,671,771 \\
cycle $E/\mathbb{F}_{p^2}$ & residue FE & 18,283,781 \\
cycle $\widehat{E}/\mathbb{F}_{q^2}$ & Weil fallback & 200,977,272 \\
cycle $\widehat{E}/\mathbb{F}_{q^2}$ & prepared-Ate raw $F_{\mathrm{num}}/F_{\mathrm{den}}$ & 76,422,749 \\
cycle $\widehat{E}/\mathbb{F}_{q^2}$ & prepared-Ate + residue & 106,441,661 \\
\midrule
Полный lollipop-cycle & stick + $E$ + улучшенный $\widehat{E}$ & 133,397,213 \\
\bottomrule
\end{tabularx}
\end{table}

Для fixed-cache режима дополнительно сравнивались calldata+commitment и
code-shards:
\[
    135{,}916{,}985\ \Gas
    \quad\text{против}\quad
    133{,}651{,}154\ \Gas.
\]
Code-shards экономят около $2.27$ млн gas за счет исключения повторной передачи
примерно $273$ КБ подготовленных линий.

Lollipop-305 показывает важный положительный факт: уменьшение числа limb-ов
действительно дает почти трехкратный выигрыш на базовой арифметике. Но этого
недостаточно для автоматического выигрыша полного цикла: структура pairing,
степень расширения и возможность удаления знаменателей не менее важны, чем
размер поля.

\section{Сводные результаты и trade-off}

\begin{longtable}{p{0.31\linewidth}p{0.36\linewidth}r}
\caption{Основные результаты работы}\\
\toprule
Подход & Что измеряется & Результат \\
\midrule
\endfirsthead
\toprule
Подход & Что измеряется & Результат \\
\midrule
\endhead
Наивный Tate baseline & строгая нижняя экстраполяция & 2.55 млрд gas \\
MNT4 full optimized on-chain & fixed-$Q$, линии строятся on-chain & 258.8 млн gas \\
MNT4 prepared sparse & кэш линий подготовлен заранее & 79.6 млн gas \\
MNT4 Article640 residue & fixed-shards equation с проверкой $G_1$ & 93.7 млн gas \\
MNT6 full equation & два Miller loop + packed FE & 226.1 млн gas \\
MNT6 Article640 residue & общий multi-Miller accumulator & 172.0 млн gas \\
KZG opening BN254 & короткая on-chain opening-проверка & 133 тыс. gas \\
KZG + MNT4 relation & BN254 non-native sparse Miller estimate & 63.3 млн constraints \\
Merkle/ordinary-FRI & ожидаемая модельная стоимость & 78.3 млн gas \\
Merkle/ordinary-FRI & размер доказательства & 2.07 МБ \\
MNT-cycle-native relation & MNT4 + MNT6 fragments & 73,068 операций \\
lollipop-305 & полный улучшенный исследовательский цикл & 133.4 млн gas \\
\bottomrule
\end{longtable}

\subsection{Три формы стоимости}

Полученные результаты показывают не один, а три взаимосвязанных вида расходов.

\paragraph{On-chain arithmetic.}
Если считать MNT-сопряжение непосредственно в EVM, приходится платить gas за
трехсловную арифметику, расширения поля и сотни шагов цикла Миллера.

\paragraph{Off-chain constraints.}
Если перенести вычисление в proof над BN254, MNT-арифметика становится
non-native. Gas on-chain уменьшается, но растет число constraints, размер
witness и время работы prover-а.

\paragraph{Calldata и proof size.}
Если использовать нативные таблицы и Merkle/FRI, можно уменьшить зависимость от
BN254 non-native arithmetic, но стоимость возвращается через Merkle-пути,
раскрытия таблиц и calldata.

\subsection{Экономическая интерпретация}

Для иллюстрации использован снимок на
\textbf{1 июня 2026 года, 13:51 MSK}. На момент снимка Etherscan
показывал:
\[
    \text{ETH price}\approx1982.85\text{ USD},
\]
\[
    \text{Ethereum standard gas price}\approx0.294\text{ gwei}.
\]
RPC Arbitrum One возвращал:
\[
    \text{Arbitrum One execution gas price}\approx0.020\text{ gwei}.
\]
Источники снимка: Etherscan Gas Tracker~\cite{etherscanGas} и публичный RPC
Arbitrum One~\cite{arbitrumRpc}. Для Ethereum mainnet:
\[
    \text{cost}_{USD}=
    \Gas\cdot \text{gas price}_{gwei}\cdot10^{-9}\cdot1982.85.
\]

\begin{table}[htbp]
\centering
\caption{Стоимость эквивалентного EVM-исполнения по снимку 01.06.2026}
\small
\begin{tabularx}{\linewidth}{Y r r r}
\toprule
Режим & Gas & Ethereum mainnet & Arbitrum L2 execution \\
\midrule
MNT4 prepared sparse & 79,586,596 & \$46.40 & \$3.16 \\
MNT4 Article640 residue & 93,734,789 & \$54.64 & \$3.72 \\
MNT6 Article640 residue & 172,004,717 & \$100.27 & \$6.82 \\
lollipop full research-cycle & 133,397,213 & \$77.76 & \$5.29 \\
Merkle/FRI lower bound & 51,359,352 & \$29.94 & \$2.04 \\
Merkle/FRI expected model & 78,340,624 & \$45.67 & \$3.11 \\
\bottomrule
\end{tabularx}
\end{table}

Столбец Arbitrum показывает только стоимость L2-исполнения одинакового числа
EVM gas. Реальная комиссия Arbitrum дополнительно включает стоимость
публикации данных в родительскую сеть. Согласно документации Arbitrum,
пользователь платит за L2 execution и за L1 data posting~\cite{arbitrumFees}.
Для prepared blob и Merkle/FRI пути эта оговорка принципиальна: большой
calldata увеличивает именно L1 data fee. Кроме того, очень крупный монолитный
вызов может потребовать разбиения на несколько транзакций даже в L2.

Таким образом, L2 снижает денежную цену арифметики, но не устраняет
архитектурное ограничение. Варианты с большим calldata остаются дорогими по
публикации данных, а тяжелые on-chain варианты требуют много ресурсов
исполнения и не всегда помещаются в один вызов.

\section{Итоговые выводы}

\subsection{Закрытие замечаний научного руководителя}

\begin{longtable}{p{0.28\linewidth}p{0.63\linewidth}}
\caption{Соответствие выполненной работы замечаниям руководителя}\\
\toprule
Замечание & Что выполнено \\
\midrule
\endfirsthead
\toprule
Замечание & Что выполнено \\
\midrule
\endhead
Сравнить алгоритмы умножения и редукции &
Реализованы и измерены Montgomery/CIOS, Barrett, FIOS, Comba/SOS,
square-Comba, lazy reduction и branchless reduction. Для башен расширений
исследованы Karatsuba и специальные умножения на нерезиденты. \\
\addlinespace
Обосновать оптимальность реализации &
Построена opcode-only нижняя граница 3-limb умножения, реализован ручной
all-stack Yul hot path, исследованы pointer/packed API, scratch arena,
fused sparse multiplication и потоковое чтение code-shards. После fused
оптимизации локальный резерв улучшения составляет доли процента полного
Miller path. \\
\addlinespace
Дать оценки сложности &
Оценки построены на нескольких уровнях: EVM opcode, $F_q$,
$F_{q^2}/F_{q^4}$ и $F_{q^3}/F_{q^6}$ операции, число Miller rounds,
стоимость FE, calldata, Merkle paths и constraints. Каждая интегральная
цифра сопровождается указанием: измерение, исполняемая модель или ручная
оценка. \\
\addlinespace
Исследовать более совершенные техники &
Реализованы shared multi-Miller accumulator, prepared sparse lines,
residue-проверка FE, KZG-opening, Merkle/DEEP-FRI отрицательный эксперимент,
ordinary-FRI cost model, MNT4/MNT6 cycle-native accounting и
lollipop-305 по ePrint 2024/1627. \\
\addlinespace
Провести всестороннее исследование &
Сравнены три фундаментальные формы стоимости: on-chain arithmetic,
off-chain constraints и calldata/proof size. Получен граничный результат:
локальные Solidity-оптимизации существенно помогают, но не устраняют
системный trade-off без нового proof/folding слоя. \\
\bottomrule
\end{longtable}

\begin{enumerate}
    \item Реализована и измерена оптимизированная 753-битная арифметика
    MNT4-753/MNT6-753 в Solidity/Yul. Для MNT4 экспериментально обоснован выбор
    Montgomery/CIOS, Karatsuba, специализированных умножений на нерезиденты,
    sparse lines и packed hot path.

    \item Построена оптимизационная лестница от наивного Tate baseline
    ($2.55$ млрд gas по строгой нижней экстраполяции) до prepared MNT4 пути
    ($79.6$ млн gas). Это показывает эффект оптимизаций и одновременно
    сохраняющуюся границу EVM.

    \item Реализована идея residue-проверки по ePrint 2024/640. Для MNT4
    финальная часть сокращается с $28.25$ млн до $6.13$ млн gas overhead.
    Для MNT6 общий multi-Miller residue verifier снижает equation baseline на
    $54.07$ млн gas.

    \item Исследована замена цикла Миллера полиномиальной проверкой. KZG
    уменьшает on-chain gas, но переносит стоимость в BN254 non-native
    constraints. Merkle/FRI избегает части этого переноса, но требует большого
    proof и calldata.

    \item Реализован предварительный MNT4/MNT6 cycle-native слой. Его
    accounting дает $24{,}126$ и $48{,}942$ multiplication operations для
    подготовленных fragments. Это не готовый CycleFold, но обоснованный
    строительный блок будущего folding.

    \item Исследована конструкция lollipop-305. Двухсловная арифметика
    действительно дешевле MNT4 примерно в $2.8$--$3.3$ раза. Однако третья
    кривая требует двух аккумуляторов числителя и знаменателя, поэтому полный
    прямой EVM-цикл остается дорогим.
\end{enumerate}

Главный вывод работы:
\begin{quote}
Без MNT-precompile или нового proof-system слоя нельзя одновременно получить
малую on-chain стоимость и малую off-chain стоимость простым переносом
вычислений. Выполненная работа локализует причины этого ограничения,
подтверждает их кодом и измерениями и показывает наиболее перспективное
направление продолжения: MNT4/MNT6 cycle-native relation layer с последующим
folding.
\end{quote}

\section{Воспроизводимость}

Финальный очищенный проект расположен в каталоге:
\[
    \code{/Users/a.i.semenov/diploma-final}.
\]
Основной запуск:
\begin{verbatim}
cd /Users/a.i.semenov/diploma-final
./scripts/run_all.sh
\end{verbatim}
Он выполняет gas-бенчмарки арифметики, MNT4/MNT6 Article640-модули, Rust
fixture cross-check, MNT-cycle accounting и ordinary-FRI cost model.

Полный аудит завершен успешно:
\begin{itemize}
    \item MNT6 Article640: $13/13$ Foundry-тестов;
    \item Rust fixture MNT6 побайтно совпадает с tracked JSON;
    \item MNT-cycle accounting: $5/5$ тестов;
    \item ordinary-FRI cost model воспроизводится;
    \item основной MNT6 verifier имеет runtime-размер $20{,}312$ байт, что
    укладывается в EIP-170.
\end{itemize}

\begin{thebibliography}{9}
\bibitem{eprint640}
Andrija Novakovic, Liam Eagen.
\emph{On Proving Pairings}.
Cryptology ePrint Archive, Paper 2024/640, 2024.
URL: \url{https://eprint.iacr.org/2024/640.pdf}.

\bibitem{eprint1627}
Craig Costello, Pratyush Korpal.
\emph{Cycles of supersingular elliptic curves for pairing-based proof systems}.
Cryptology ePrint Archive, Paper 2024/1627, 2024.
URL: \url{https://eprint.iacr.org/2024/1627.pdf}.

\bibitem{sonobe}
Privacy and Scaling Explorations.
\emph{Sonobe: a modular library for folding schemes}.
URL: \url{https://github.com/privacy-ethereum/sonobe}.

\bibitem{etherscanGas}
Etherscan.
\emph{Ethereum Gas Tracker}.
URL: \url{https://etherscan.io/gastracker}.
Снимок данных: 01.06.2026, 13:51 MSK.

\bibitem{arbitrumRpc}
Offchain Labs.
\emph{Arbitrum One public RPC endpoint}.
URL: \url{https://arb1.arbitrum.io/rpc}.
Снимок данных: 01.06.2026, 13:51 MSK.

\bibitem{arbitrumFees}
Offchain Labs.
\emph{Arbitrum Nitro: gas and fees}.
URL: \url{https://docs.arbitrum.io/how-arbitrum-works/gas-fees}.
\end{thebibliography}

\end{document}
